% Generated draft: requires mathematical and source review.
\documentclass[11pt]{article}
\usepackage[margin=25mm]{geometry}
\usepackage{fontspec}
\setmainfont{Latin Modern Roman}
\usepackage{amsmath,amssymb,mathtools,unicode-math}
\setmathfont{Latin Modern Math}
\usepackage{enumitem,needspace,xurl,hyperref}
\setlength{\parindent}{0pt}
\setlength{\parskip}{5pt}
\setlength{\emergencystretch}{3em}
\newcommand{\N}{\mathbb{N}}
\newcommand{\Z}{\mathbb{Z}}
\newcommand{\Q}{\mathbb{Q}}
\newcommand{\R}{\mathbb{R}}
\newcommand{\C}{\mathbb{C}}
\newcommand{\F}{\mathbb{F}}
\newcommand{\A}{\mathbb{A}}
\newcommand{\PP}{\mathbb{P}}
\newcommand{\E}{\mathbb{E}}
\newcommand{\eps}{\varepsilon}
\newcommand{\dd}{\,\mathrm{d}}
\begin{document}
\section*{UnsolvedMath ID: 801068}
\textbf{Source problem code:} catalogue entry
\subsection{Definitions and mathematical statement}
Let \(\mathcal H\) denote the class of Hénon maps intended in Question 20 of the supplied source. For \(f\in\mathcal H\), let \(F(f)\) denote its forward Fatou set, namely the normality locus of the family of iterates \(\{f^n:n\geq 0\}\), using the ambient space and normality convention intended by that source. A \emph{Fatou component} is a connected component of \(F(f)\). Such a component \(U\) is \emph{wandering} if
\[
 f^m(U)\cap f^n(U)=\varnothing
 \qquad\text{for every pair of integers }0\leq m<n.
\]
Here \(f^0\) is the identity and \(f^n\) is the \(n\)-fold iterate for \(n\geq1\).

\textbf{Question 20 (Berger, Bedford, 2023).} Does every quasi-hyperbolic map \(f\in\mathcal H\) have no wandering Fatou component? Equivalently, with the preceding conventions, is it true that
\[
 \forall f\in\mathcal H,\quad
 f\text{ is quasi-hyperbolic}
 \ \Longrightarrow\ 
 \forall U\in\pi_0(F(f)),\ 
 \exists\,0\leq m<n:\quad
 f^m(U)\cap f^n(U)\neq\varnothing,
\]
where \(\pi_0(F(f))\) denotes the set of connected components of \(F(f)\)?

\textbf{Unresolved qualifications.} The supplied excerpt does not specify the precise class \(\mathcal H\), its ambient space or parameter restrictions, the definition of quasi-hyperbolicity, or the normality convention defining the Fatou set. These remain source-dependent terms in this formulation; the excerpt alone does not determine a fully specified mathematical assertion.
\subsection{Short English statement}
Berger and Bedford ask whether quasi-hyperbolicity rules out wandering Fatou components for the Hénon maps under consideration. The supplied excerpt leaves the map class and several defining conventions unspecified, so this statement requires review against the source.
\subsection{Sources}
\begin{itemize}
\item Local UnsolvedMath registry, AMR Open Problem Lists: permanent ID \(801068\), exact source problem code \texttt{catalogue entry}, title \emph{Henon Maps: Open Problems --- Question 20}. The supplied registry labels the problem \texttt{open}; this is supplied status metadata, not an independently verified assessment.
\item Berger, Bedford, \emph{Henon Maps: A List of Open Problems} (2023), Question 20. Supplied source URL: \url{https://arxiv.org/abs/2312.03907}. The record attributes the question to Berger and Bedford, with proposed year 2023.
\item Supplied provenance: accessed 2026-10-16; extraction method \texttt{arxiv\_latex}; upstream revision \texttt{d371ee917ac9647fbb06cd6fc529deec6f3dda82}. These are supplied provenance details; no independent source consultation or verification was performed.
\end{itemize}
\end{document}
